\section{Introduction}

The fine-structure constant is one of the central dimensionless constants
of low-energy physics. In conventional notation,
\[
\alpha=\frac{e^2}{4\pi\varepsilon_0\hbar c},
\]
or, in naturalized electromagnetic units, it is the dimensionless
coupling governing the strength of electromagnetic interactions. It
appears in atomic binding, spectroscopy, radiative corrections,
scattering amplitudes, and precision metrology. Its inverse value is
approximately
\[
\alpha^{-1}\approx 137.035999177.
\]

In standard physics this value is not derived from a deeper structural
principle. It is inferred through operational routes such as atom recoil,
spectroscopy, electron \(g-2\), and electrical standards, and then used
as an input to low-energy quantum electrodynamics and atomic theory
\cite{NIST_CODATA_2022_alpha_inv,Parker2018_alpha_Cs,
Morel2020_alpha_Rb,Aoyama2019_electron_gminus2_review,
BIPM_SI_Brochure_9th}.

The purpose of this paper is to present a compact-fiber structural
prediction for the low-energy electromagnetic bridge coupling. The
derived object is denoted
\[
\alpha_F,
\qquad
x=\alpha_F^{-1}.
\]
The central result is that compact-fiber carrier structure,
radiative/vibrational bridge placement, cover-mediated feedback,
self-consistent transport sharing, and the first closed return yield a
quartic equation for \(x\). Its positive root is
\[
x_{\rm fiber}=137.0359991771859\ldots,
\]
so that
\[
\alpha_F=\frac{1}{x_{\rm fiber}}.
\]
The measured value of \(\alpha\) is not used as an algebraic input to
the bridge equation.

\subsection{Structural target}

The compact-fiber framework begins from the finite carrier
\[
F=Z_4\times Z_4\times Z_8,
\]
with
\[
|F|=128,
\qquad
N_m=4,
\qquad
N_e=8.
\]
The two \(Z_4\) factors represent the paired magnetic sector, while the
\(Z_8\) factor represents the electric-cover sector. The relation
\[
N_e=2N_m
\]
supplies the canonical electric-cover / magnetic-base quotient used in
the cover-mediated feedback term. The present paper does not rederive
the compact fiber from first principles; it takes the carrier and sector
capacities from the foundational compact-fiber construction and derives
the electromagnetic bridge coupling from them
\cite{Wells2026PaperI,Wells2026PaperII,Wells2026PaperIII}.

The electromagnetic coupling is treated here as a carrier-level
material/radiative bridge normalization. It is not an element-by-element
function of specific labels
\[
f\in F,
\qquad
v\in V_{\rm rad}.
\]
Instead,
\[
\alpha_F(f,v)=\alpha_F
\]
for all bridge-admissible element labels. The relevant structural
question is therefore how material carrier access and
radiative/vibrational bridge-placement access combine at the level of
the electromagnetic bridge, not how many ordered pairs exist in
\[
F\times V_{\rm rad}.
\]

\subsection{Derivation strategy}

The derivation proceeds through four structural steps. First, the
material compact-fiber carrier and the radiative/vibrational
bridge-placement carrier determine the base bridge count
\[
B=128+9=137.
\]
Second, cover-mediated feedback from the paired magnetic sector into the
electromagnetic bridge determines the feedback burden
\[
K=\frac{74}{15}.
\]
Third, self-consistent transport sharing gives the first-order
fixed-point structure
\[
x=B+\frac{K}{x}.
\]
Fourth, the first closed return of the admitted feedback supplies the
second-order return term. These ingredients yield the compact bridge
equation
\[
x
=
B+\frac{K}{x}
-
\frac{K N_m}{(N_e-1)x^3}.
\]
The detailed derivation of each term is given in the sections that
follow.

\subsection{Identification with the low-energy electromagnetic coupling}

The fixed-point derivation produces a compact-fiber structural bridge
coupling
\[
\alpha_F=\frac{1}{x_{\rm fiber}}.
\]
The identification
\[
\alpha_F=\alpha_{\rm low-energy}
\]
is made by physical role, not by numerical agreement alone. The
compact-fiber electromagnetic bridge is the material/radiative interface.
A leading atomic binding observable is a closed material-radiative
source-response process with two bridge vertices,
\[
\text{material source}
\rightarrow
\text{radiative/electromagnetic bridge}
\rightarrow
\text{material response}.
\]
Thus the leading atomic scale carries two bridge factors,
\[
E_{\rm atomic,F}
=
C_{\rm mass/readout}\alpha_F^2.
\]
This matches the role of \(\alpha^2\) in the Hartree/Rydberg scale,
where
\[
E_h=\alpha^2m_ec^2,
\qquad
R_\infty=\frac{\alpha^2m_ec}{2h}.
\]
This paper therefore treats \(\alpha_F\) as the compact-fiber structural
counterpart of the low-energy fine-structure constant.

\subsection{Scope of the paper}

The paper separates the structural coupling derivation from operational
measurement routes. It derives the compact-fiber bridge coupling
\(\alpha_F\), identifies it with the low-energy electromagnetic coupling
by its \(\alpha_F^2\) role in leading atomic binding, and gives a
controlled correspondence to empirical extraction routes.

The result is a conditional structural derivation: given the named
compact-fiber framework lemmas, the fixed-point equations and their roots
follow algebraically. It is not presented as an unconditional derivation
from bare postulates alone. The paper also does not claim a completed
route-by-route derivation of CODATA or metrological extraction methods;
full route-specific metrological closure remains downstream work.

The measured value of \(\alpha\) is used only after the derivation, as
an empirical consistency check. No measured value of \(\alpha\) is used
to set \(B\), \(K\), the closed-return coefficient, or any continuously
adjusted coefficient in the bridge equation.